\renewcommand{\sectionspath}[1]{tox_manual/fortran_pipeline/#1}

\lstset{ language=Fortran, basicstyle=\ttfamily\small, keywordstyle=\color{blue}, commentstyle=\color{green!40!black}, stringstyle=\color{purple}, showstringspaces=false, breaklines=true, frame=single, rulecolor=\color{black!30}, backgroundcolor=\color{black!5} }

\subsection{Fortran Pipeline}
\subsubsection{Overview}
\begin{enumerate}
    \item Data dependencies and relationships All arrays read have dependencies to
          each other. The gene ids array stores the ids, column number seven of the expression
          vectors belongs to entry number seven in the gene ids. Entry number nine
          in the gene to family mapping array, belongs to entry nine in the gene ids.
          These dependencies must be given at any time, therefore, only the provided
          functions shall be used to modify arrays in any form. For error code details
          please check \hyperref[sec:error-handling]{error handling section}.
          
    \item Error handling reference (placeholder for now)
          
    \item General usage principle (Placeholder for now)
\end{enumerate}

\input{\sectionspath{data_management}}
\input{\sectionspath{data_persistence}}
\input{\sectionspath{data_processing}}
\input{\sectionspath{family_based_analysis}}
\input{\sectionspath{rap_projection_and_analysis_functions}}
\input{\sectionspath{paralog_analysis}}
\input{\sectionspath{trajectory_contribution_analysis}}
\input{\sectionspath{fortran_utils}}